\section{Summary, engineering consequences and open problems}
\label{concl:sec}

\begin{table}[ht]
\centering\footnotesize
\renewcommand{\arraystretch}{1.25}
\begin{tabularx}{\textwidth}{@{}>{\raggedright\arraybackslash}p{0.17\textwidth}>{\raggedright\arraybackslash}p{0.17\textwidth}>{\raggedright\arraybackslash}p{0.17\textwidth}>{\raggedright\arraybackslash}p{0.14\textwidth}>{\raggedright\arraybackslash}X@{}}
\toprule
Model (application) & Singularity & Non-negative modes & Gap $\omega_0$ & Status in these notes\\
\midrule
Hyperbolic MEMS $w^{-q}$, $q>1$ (electrostatic, vdW, Casimir actuators)
 & touchdown $w\simeq c(T-t)^{\frac2{q+1}}$, flat in cone
 & $1$ (time), $0$ (boosts)
 & $\min\{1,\frac{2(q-1)}{q+1}\}$
 & spectrum proved; stability modulo symmetries and with damping sketched
   (\cref{mems:thm-main,mems:prop-damped}); new\\
MEMS, focused profiles $\psi_n$, $d=2,3$
 & touchdown at a point, $\psi_n(0)\to0$
 & $n$ extra unstable
 & --
 & numerical; \cref{mems:conj-prof}\\
Reduced hyperbolic Gelfand $u_{tt}-\Delta u=\delta e^u$ (Cattaneo heat, FK source)
 & $u\simeq\log\frac2\delta-2\log(T-t)$
 & $1$, $0$
 & $1$ (sharp in $d=1$)
 & stability sketched (\cref{therm:thm:reduced}); 1D exact check via Liouville;
   cf.\ \cite{Kichenassamy1996}\\
Full Cattaneo--FK model
 & $\theta\simeq-\log(T-t)$ (parabolic rate)
 & $1$, $0$ (tilts)
 & $1$ (Jordan block at $-1$)
 & sketched, less checked (\cref{therm:thm:full})\\
Arrhenius source
 & no blowup
 & --
 & --
 & FK blowup only as an intermediate regime\\
Softening foundation / $\varphi^4$ / NLKG, $|u|^{p-1}u$
 & $u\simeq\kappa(T-t)^{-\frac2{p-1}}$
 & $1$, $0$
 & $1$
 & known for $\gamma=m=0$ \cite{Ostermann2024}; lower-order lemma sketched
   (\cref{power:thm-lower-order})\\
Co-rotational wave maps (inertial director dynamics)
 & $2\arctan\frac{r}{\sqrt{d-2}(T-t)}$
 & $1$
 & $\approx0.54$ ($d=3$)
 & known (\cref{geom:thm:d3}); spectrum reproduced numerically\\
Relativistic membranes / Born--Infeld
 & shrinking spheres (lightlike collapse)
 & infinitely many
 & --
 & Donninger method not applicable to spheres; general case open
   (\cref{geom:sec:membranes})\\
Acoustic toy model $u_{tt}-\Delta u=u_t^2$
 & $u\simeq-\log(T-t)$
 & $1$, $0$ (shift; tilts)
 & $1$
 & radial sketched (\cref{acou:thm-rad}); nonradial conjectural; cf.\ \cite{Speck2020}\\
Inviscid Westervelt
 & degeneration $1-2kp\propto(T-t)^{1/2}$
 & $1$, $\frac14$ (tilt), $0$
 & formal
 & quasilinear, open (\cref{acou:sec})\\
\bottomrule
\end{tabularx}
\caption{Status of the models. ``Non-negative modes'' lists the eigenvalues
$\lambda\ge0$ of the linearisation in similarity variables together with
their origin. ``Sketched'' means a proof sketch with all key estimates in
the style of \cite{DonningerSchoerkhuber2017,Ostermann2024}.}
\label{concl:tab}
\end{table}

\subsection{What the analysis says about the devices}
\begin{enumerate}[leftmargin=1.5em]
\item \emph{Rates are universal and robust.} Once inertia is retained, the
final approach to touchdown or runaway is governed by the spatially
homogeneous ODE. For the Coulomb MEMS force the gap closes like
$(9\lambda/2)^{1/3}(T-t)^{2/3}$, against $(3\lambda(T-t))^{1/3}$ in the
parabolic (overdamped) model. Damping, restoring forces and smooth variations
of the permittivity change only lower-order corrections. They decay in
similarity variables, so the hyperbolic rate is the true terminal one even
for strongly damped devices. The parabolic rate is at most an intermediate
regime.
\item \emph{The terminal profile is flat.} In the backward light cone the
stable profile is constant (or a Lorentz-tilted constant). Touchdown
therefore occurs along a spacelike hypersurface $t=T^*(x)$, not at an
isolated point as in the parabolic theory. Focused, point-like touchdown
profiles exist for $d=2,3$ but are unstable. A generic device is not expected
to show them.
\item \emph{Boundaries do not matter locally.} The results are local in the
light cone of the singular point. Clamping conditions, the geometry of the
device and far-away material parameters enter only through the touchdown
time $T^*$ and the tilt $a^*$.
\item \emph{Which singularity actually forms is model-dependent.} The full
Maxwell--Cattaneo thermal-explosion model blows up at the \emph{parabolic}
logarithmic rate, not at the hyperbolic one of the frequently used reduced
model. Arrhenius kinetics removes the blowup altogether. Viscous nonlinear
acoustics leaves the light-cone setting entirely.
\end{enumerate}

\subsection{Open problems suggested by these notes}
\begin{enumerate}[leftmargin=1.5em]
\item Turn the proof sketches of
\cref{mems:thm-main,mems:prop-damped,therm:thm:reduced} into complete
proofs. The new ingredients relative to \cite{Ostermann2024} are the
negative scaling exponent (finite-rank splitting of the free operator,
\cref{mems:lem-free}) and the need to stay away from $w=0$, which the
$L^\infty$ control gives.
\item Prove the existence of the focused MEMS profiles $\psi_n$ in $d=2,3$
(\cref{mems:conj-prof}) and determine their co-dimension. A shooting argument
around the singular static solution $A|x|^{2/(q+1)}$, as in
\cite{BizonMaisonWasserman2007}, looks feasible.
\item Go beyond the light cone: regularity of the touchdown hypersurface,
$a^*=\nabla T^*$, and the hyperboloidal approach of
\cite{BiernatDonningerSchoerkhuber2021,ChenEtAl2024} for MEMS.
\item Models not covered by the method in its present form: MEMS plates
($\Delta^2$, no light cone), nonlocal MEMS, the full Cattaneo model in
higher regularity, and the quasilinear problems (Westervelt degeneration,
membranes).
\item Compare the rigorous picture of \cref{mems:sec} in detail with the
formal analysis and numerics of \cite{GimperleinHeLacey2022,He2021}. This
could not be done here because the paper was not accessible.
\end{enumerate}
